\begin{textsample}{POS Dim 2 – 2017\_c6 – Score -1.00 – 2017\_c6/enade\_2017\_eng\_civ\_9cc7e908b9.txt}  \label{ex:f2_pos_272}
CAPÍTULO 6 ANÁLISE TÉCNICA DA PROVA Este capítulo tem por objetivo apresentar o desempenho dos estudantes concluintes de Engenharia Civil no Enade/2017.
Para isso, foram calculadas as estatísticas básicas da prova em seu todo ( seção 6.1.1 ), bem como as estatísticas relacionadas aos Componentes de Formação Geral ( seção 6.1.2 ) e de Conhecimento Específico da Área ( seção 6.1.3 ).
Dadas as suas características, foram analisadas, em separado, as questões objetivas ( seção 6.2 ) e as questões discursivas ( seção 6.3 ).
Para as questões objetivas, foram disponibilizados os índices de facilidade e de discriminação ponto-bisserial, também em separado, para os Componentes de Formação Geral ( seção 6.2.1 ) e de Conhecimento Específico ( seção 6.2.2 ).
De cada componente, uma das questões foi escolhida para exemplificar a análise gráfica, relacionando as alternativas escolhidas pelos estudantes ( inclusive o gabarito ) com o número de acertos no componente.
O Anexo I apresenta a íntegra da análise gráfica para todas as questões objetivas.
Para cada uma das questões discursivas, os conteúdos dos tipos mais comuns de respostas dos estudantes são apresentados e comparados com o padrão de respostas esperado ( ver Anexo VIII com o padrão de respostas ).
Tomando como base as duas questões discursivas do Componente de Formação Geral, a seção 6.3.1.6 apresenta comentários sobre a correção das respostas com respeito à Língua Portuguesa.
As tabelas apresentam as seguintes estatísticas das notas19 : média do desempenho na prova, erro padrão da média, desvio padrão, nota mínima, mediana e nota máxima para cada um de seus componentes.
Tais estatísticas contemplam o total de estudantes concluintes da Área de Engenharia Civil inscritos e presentes à prova do Enade/2017, tendo em vista agregações, ou por Grandes Regiões e o país como um todo, ou por Categoria Administrativa, Organização Acadêmica da IES e Modalidade de Ensino.
Em relação aos gráficos de distribuição de notas, o intervalo considerado foi de 10 unidades, aberto à esquerda e fechado à direita, com exceção do primeiro intervalo, [ 0 ; 10 ], fechado em ambos os extremos.
Para os gráficos de distribuição das notas das questões discursivas, foram consideradas mais duas categorias : questão em branco20 e nota zero. 19 Uma definição dessas estatísticas pode ser encontrada no Glossário. 20 Nesse grupo estão incluídas também as respostas classificadas como nulas ou desconsideradas. 6.1 ESTATÍSTICAS BÁSICAS DA PROVA Esta seção apresenta estatísticas selecionadas e histogramas da nota geral ( 6.1.1 ) e de cada componente : Formação Geral ( 6.1.2 ) e Conhecimento Específico ( 6.1.3 ).
São também apresentadas estatísticas selecionadas de subpopulações caracterizadas por Grande Região, Categoria Administrativa, Organização Acadêmica e Modalidade de Ensino. 6.1.1 Estatísticas Básicas Gerais A Tabela 6.1 apresenta as Estatísticas Básicas da prova, por Grande Região, dos estudantes concluintes de Engenharia Civil.
A Média das notas da prova como um todo ( nas seções seguintes serão analisados os Componentes de Formação Geral e de Conhecimento Específico ) foi 44,1, sendo que os alunos da região Sul obtiveram a Média mais alta ( 48,0 ), e os da região Norte obtiveram a Média mais baixa ( 41,6 ).
As demais médias foram : 45,0 na região Nordeste, 42,7 na região Sudeste e 44,5 na região Centro-Oeste.
O Desvio padrão para o Brasil como um todo foi 13,4, sendo o maior Desvio padrão encontrado na região Nordeste ( 13,8 ) e o menor, na região Norte ( 12,8 ), indicando uma dispersão um pouco menor das notas desta última região.
A região que obteve a maior nota Máxima foi a Sudeste ( 96,4 ), ao passo que a região que atingiu a menor nota Máxima foi a Centro-Oeste ( 85,6 ).
A Mediana do Brasil como um todo foi 43,4, sendo a maior Mediana obtida na região Sul ( 47,7 ), e a menor obtida na região Norte ( 40,9 ).
A nota Mínima foi zero em todas as regiões, exceto na região Norte ( 5,1 ).
Considerando -se as notas segundo Grande Região, observa -se que existe diferença estatisticamente significativa ao nível de 95 \% entre a menor Média, obtida na região Norte ( 41,6 ), e as médias de todas as demais regiões21. 21 Todas as comparações deste capítulo utilizam os intervalos de 95 \%.
Os erros-padrão da média que possibilitam os testes estão disponíveis nas tabelas.
O comportamento das notas dos estudantes de todo o Brasil pode ser observado no Gráfico 6.1 que apresenta um histograma com a distribuição das mesmas.
Essa é uma distribuição unimodal com a moda no intervalo ( 40 ; 50 ], seguido pelo intervalo ( 30 ; 40 ].
A Tabela 6.2 apresenta informações referentes à Média da nota final desagregadas por Categoria Administrativa, por Organização Acadêmica e por Modalidade de Ensino.
A Média das notas da prova como um todo foi 44,1.
Em relação à Categoria Administrativa, os estudantes das IES Públicas obtiveram Média mais alta ( 54,1 ), e os das IES Privadas Média mais baixa ( 41,8 ), que a Média nacional.
Observa -se que existe diferença estatisticamente significativa entre as médias das notas das IES Públicas e das Privadas.
A diferença entre as médias das regiões Norte e Sul ( 6,4 ), a menor e a maior Média, é inferior à diferença entre IES Públicas e Privadas ( 12,3 ), caracterizando uma maior diversidade administrativa do que regional.
O Desvio padrão para as IES Públicas ( 14,0 ) foi superior ao do Brasil como um todo ( 13,4 ), indicando uma dispersão um pouco maior das notas nessa Categoria Administrativa.
No tocante à Organização Acadêmica, as Universidades e os CEFET/IFET obtiveram Média mais alta que a nacional ( 46,2 e 53,9, respectivamente ).
As Médias de Faculdades e de Centros universitários foram menores que a nacional, 40,0 e 42,4, respectivamente.
Constata -se que existe diferença estatisticamente significativa ao nível de 95 \% entre a Média mais baixa, 40,0 das Faculdades, e a dos outros tipos de Organização Acadêmica.
A Média da Modalidade de Educação a Distância ( 40,7 ) foi inferior à Média da Educação Presencial ( 44,1 ), e a diferença entre elas é estatisticamente significativa ao nível de 95 \%. 6.1.2 Estatísticas Básicas no Componente de Formação Geral A Tabela 6.3 apresenta as Estatísticas Básicas em relação ao componente da prova que avalia a Formação Geral dos estudantes concluintes.
Os alunos de todo o Brasil obtiveram desempenho médio de 54,7.
Quanto à variabilidade, o Desvio padrão das notas dos estudantes do Brasil, como um todo, foi 17,6.
A maior Média foi obtida na região Sul ( 57,7 ), e a menor, na região Norte ( 53,1 ).
As demais médias foram : 55,5 na região Nordeste, 53,7 na região Sudeste e 54,9 na região Centro-Oeste.
Já o maior Desvio padrão foi obtido na região Norte ( 18,0 ), e o menor, na região Sul ( 16,9 ).
Os demais desvios padrões foram : 17,7 na região Nordeste, 17,5 na região Sudeste e 17,9 na região Centro-Oeste.
A maior nota no Componente de Formação Geral da prova do Enade/2017 foi 99,2, obtida por, pelo menos, um aluno na região Sudeste.
A menor nota Máxima foi obtida na região Sul ( 98,0 ).
A Mediana do Brasil, como um todo, foi 55,8, sendo a menor Mediana encontrada na região Norte ( 53,8 ), e a maior, encontrada na região Sul ( 59,1 ).
A nota Mínima nessa parte foi zero em todas as regiões, sem exceção.
Considerando -se as notas segundo Grande Região, observa -se que existe diferença estatisticamente significativa entre a maior Média das notas do Componente de Formação Geral, obtida na região Sul ( 57,7 ), e as médias obtidas nas demais regiões.
O Gráfico 6.2 propicia a avaliação do desempenho dos estudantes do Componente de Formação Geral, a partir do histograma da distribuição das notas correspondentes.
A distribuição é unimodal, com moda em ( 50 ; 60 ], um intervalo acima do intervalo modal da distribuição de notas da prova, como um todo ( Gráfico 6.1 ).
Nota -se, ainda que, no Gráfico 6.2, as notas apresentam dispersão um pouco maior do que a do Gráfico 6.1 ( distribuição das notas da prova ), confirmado pela comparação dos desvios padrões : 13,4 para a nota da prova como um todo e 17,6 para o Componente de Formação Geral.
Na Tabela 6.4 são apresentadas as informações referentes ao desempenho dos concluintes do Componente de Formação Geral, em diferentes agregações : Categoria Administrativa, Organização Acadêmica e Modalidade de Ensino.
Observa -se que existe diferença estatisticamente significativa entre as médias dos tipos de Categoria Administrativa : a maior Média obtida por estudantes de IES Públicas ( 63,5 ) e a menor pelos de IES Privadas ( 52,7 ).
Nota -se que também há diferença estatisticamente significativa entre as médias dos CEFET/IFET ( 65,6 ), a maior, e de todas as demais categorias : Universidades ( 56,7 ), Centros Universitários ( 53,1 ) e Faculdades ( 50,9 ), a menor.
No que se refere à Modalidade de Ensino, a Educação Presencial apresentou maior Média do que a Educação a Distância ( médias 54,7 e 52,5, respectivamente ).
Tais médias apresentam diferenças estatisticamente significativas. 6.1.3 Estatísticas Básicas do Componente de Conhecimento Específico A Tabela 6.5 apresenta as Estatísticas Básicas referentes ao Componente de Conhecimento Específico da área de Engenharia Civil.
A Média do desempenho dos alunos do Brasil, como um todo, foi 40,5.
A maior Média foi obtida na região Sul ( 44,7 ), e a menor, na região Norte ( 37,8 ).
As demais médias foram : 41,5 na região Nordeste, 39,1 na região Sudeste e 41,0 na região Centro-Oeste.
Quanto à variabilidade das notas, o Desvio padrão do Brasil, como um todo, foi 14,6, sendo o maior Desvio padrão observado na região Nordeste ( 15,1 ), e o menor, na região Norte ( 13,8 ).
Os demais desvios foram : 14,3 na região Sudeste, 14,4 na região Sul, e 14,6, na região Centro-Oeste.
A Mediana das notas dos estudantes de todo o Brasil foi 40,0.
A maior Mediana ocorreu na região Sul ( 44,5 ), e a menor, na região Norte ( 36,3 ).
As demais medianas foram : 40,3 na região Nordeste, 38,3 na região Sudeste e 40,3 na região Centro-Oeste.
A nota Máxima do Brasil, como um todo, foi 99,0, sendo obtida por, pelo menos, um aluno na região Sudeste.
As demais notas máximas foram : 93,7 na região Norte, 90,5 na região Nordeste, 92,5 na região Sul e 87,3 na região Centro-Oeste.
A nota Mínima foi zero em todas as regiões.
Observa -se que existe diferença estatisticamente significativa da menor Média das notas do Componente de Conhecimento Específico da região Norte ( 37,8 ) em relação às demais regiões.
Assim como os Gráficos 6.1 e 6.2, o Gráfico 6.3, apresentado a seguir, permite uma avaliação do desempenho de concluintes de Engenharia Civil em relação ao Componente de Conhecimento Específico com um histograma da distribuição das notas correspondentes.
Essa também é uma distribuição unimodal, e o grupo modal é o ( 30 ; 40 ], abaixo do grupo modal da prova como um todo e do grupo modal para a Formação Geral.
A Tabela 6.6 apresenta uma comparação dos resultados em relação à Categoria Administrativa, à Organização Acadêmica e a Modalidade de Ensino, agora levando -se em conta o desempenho de alunos do Componente de Conhecimento Específico da prova.
No que se refere à Organização Acadêmica, a maior Média foi dos CEFET/IFET ( 50,0 ), vindo a seguir a das Universidades ( 42,6 ), a dos Centros Universitários ( 38,8 ) e, depois, a das Faculdades ( 36,3 ).
O maior Desvio padrão, e acima do valor para o Brasil, como um todo ( 14,6 ), foi o das Universidades ( 15,3 ).
As Universidades também obtiveram a maior nota Máxima ( 99,0 ).
Os Centros Universitários obtiveram nota Máxima igual a 90,5, as Faculdades obtiveram nota Máxima igual a 90,0, e os CEFET/IFET igual a 84,7.
As medianas foram : 35,8 nas Faculdades, a menor delas, 38,3 nos Centros Universitários, 41,8 nas Universidades e 50,0 nos CEFET/IFET, a maior.
A nota Mínima foi zero para todas as Organizações Acadêmicas.
Observa -se que existem diferenças estatisticamente significativas ao nível de 95 \% no Componente de Conhecimento Específico entre as notas das Organizações Acadêmicas.
Quanto à Categoria Administrativa, observa -se um comportamento semelhante àquele da parte de Formação Geral e à prova como um todo, ou seja, existe diferença estatisticamente significativa entre as médias das IES Públicas ( 51,0 ) e as das IES Privadas ( 38,1 ).
Nesse caso, também, a maior Média foi obtida por alunos de IES Públicas de ensino.
Também semelhante é o comportamento das médias no que se refere toca à Modalidade de Ensino, uma vez que há diferença estatisticamente significativa entre a Média da Educação Presencial ( 40,5 ) e a da Educação a Distância ( 36,8 ). 6.2 ANÁLISE DAS QUESTÕES OBJETIVAS Esta seção apresenta estatísticas selecionadas e histogramas das Questões Objetivas de Formação Geral ( 6.2.1 ) e de Conhecimento Específico ( 6.2.2 ).
São também apresentadas e comparadas as médias das subpopulações caracterizadas por Grande Região. 6.2.1 Componente de Formação Geral A Tabela 6.7 apresenta as Estatísticas Básicas relativas às oito questões objetivas do componente da prova que abrange a Formação Geral dos estudantes.
A Média do Brasil foi 55,2.
A maior Média foi encontrada na região Sul ( 59,5 ), e a menor, nas regiões Norte e Sudeste ( 53,5 ).
As demais médias foram : 56,6 na região Nordeste e 56,5 na região Centro-Oeste.
O Desvio padrão do Brasil foi 21,6, sendo o maior Desvio padrão encontrado na região Norte ( 22,0 ), e o menor, na região Sul ( 20,8 ).
Os demais desvios foram : 21,9 na região Nordeste, 21,5 na região Sudeste e 21,8 na região Centro-Oeste.
A Mediana do Brasil como um todo e as das regiões Norte e Sudeste foram 50,0, enquanto nas demais regiões a Mediana foi 62,5.
A nota Máxima, 100,0, e a nota Mínima, zero, foram iguais para todas as regiões.
A Tabela 6.8 apresenta o Índice de Facilidade e o Índice de Discriminação ( ponto-bisserial ) para cada uma das questões objetivas do Componente de Formação Geral.
Quanto ao Índice de Facilidade, foram usadas as seguintes cores para diferenciar o nível de dificuldade da questão :  Azul para as questões classificadas com índice Muito fácil ( \&gt ; =0,86 ), verde para as questões classificadas com índice Fácil ( 0,61 a 0,85 ), amarelo para as questões classificadas com Médio ( 0,41 a 0,60 ), vermelho para as questões classificadas com Difícil ( 0,16 a 0,40 ) e roxo para as questões classificadas com Muito difícil ( \&lt ; =0,15 ).
Já quanto ao Índice de Discriminação, foram usadas as seguintes cores para qualificar a questão :  As questões classificadas com índice Fraco receberam a cor vermelha ( \&lt ; =0,19 ), as classificadas com Médio receberam a cor amarela ( 0,20 a 0,29 ), as classificadas com Bom receberam a cor verde ( 0,30 a 0,39 ) e as classificadas com Muito bom ( \&gt ; =0,40 ) receberam a cor azul.
As questões objetivas do Componente de Formação Geral, segundo o Índice de Facilidade, foram assim avaliadas : das oito questões, uma teve o Índice de Facilidade classificado como Muito fácil, e duas questões foram tidas como Fácil.
Quatro questões foram consideradas com índice de dificuldade Médio, situando -se no intervalo entre 0,41 e 0,60 do Índice de Facilidade, ou seja, houve entre 41,0 \% e 60,0 \% de acertos, enquanto uma questão foi classificada na categoria Difícil, situando -se no intervalo entre 0,16 e 0,40.
Por fim, A Mediana do Brasil como um todo e as das regiões Norte e Sudeste foram 50,0, enquanto nas demais regiões a Mediana foi 62,5.
A nota Máxima, 100,0, e a nota Mínima, zero, foram iguais para todas as regiões.
A Tabela 6.8 apresenta o Índice de Facilidade e o Índice de Discriminação ( ponto-bisserial ) para cada uma das questões objetivas do Componente de Formação Geral.
Quanto ao Índice de Facilidade, foram usadas as seguintes cores para diferenciar o nível de dificuldade da questão : • Azul para as questões classificadas com índice Muito fácil ( \&gt ; =0,86 ), verde para as questões classificadas com índice Fácil ( 0,61 a 0,85 ), amarelo para as questões classificadas com Médio ( 0,41 a 0,60 ), vermelho para as questões classificadas com Difícil ( 0,16 a 0,40 ) e roxo para as questões classificadas com Muito difícil ( \&lt ; =0,15 ).
Já quanto ao Índice de Discriminação, foram usadas as seguintes cores para qualificar a questão : • As questões classificadas com índice Fraco receberam a cor vermelha ( \&lt ; =0,19 ), as classificadas com Médio receberam a cor amarela ( 0,20 a 0,29 ), as classificadas com Bom receberam a cor verde ( 0,30 a 0,39 ) e as classificadas com Muito bom ( \&gt ; =0,40 ) receberam a cor azul.
As questões objetivas do Componente de Formação Geral, segundo o Índice de Facilidade, foram assim avaliadas : das oito questões, uma teve o Índice de Facilidade classificado como Muito fácil, e duas questões foram tidas como Fácil.
Quatro questões foram consideradas com índice de dificuldade Médio, situando -se no intervalo entre 0,41 e 0,60 do Índice de Facilidade, ou seja, houve entre 41,0 \% e 60,0 \% de acertos, enquanto uma questão foi classificada na categoria Difícil, situando -se no intervalo entre 0,16 e 0,40.
Por fim, A Tabela 6.9 apresenta a distribuição das questões considerando -se simultaneamente a classificação dos dois índices.
As oito questões com Índice de Discriminação Muito bom figuraram entre quatro níveis de dificuldade Muito fácil, Fácil, Médio e Difícil : uma classificada na categoria Muito Fácil ( a questão 6 ) do Índice de Facilidade, duas classificadas na categoria Fácil ( questões 4 e 5 ), quatro na categoria Médio ( questões 2, 3, 7 e 8 ) e, por fim, uma na categoria Difícil ( questão 1 ).
Em particular, a questão 4 foi a que apresentou o maior poder discriminatório, com índice 0,51, tendo sido considerada Fácil em termos de facilidade, com uma proporção de 0,65 acertos.
O máximo de acertos foi alcançado pela questão 6, com um Índice de Facilidade de 0,86.
As questões 4 e 5 foram consideradas Fácil, com índices de facilidade 0,65 e 0,68, respectivamente.
Já os índices de discriminação foram Muito bom para as questões 2 ( 0,41 ), 3 ( 0,45 ), 7 ( 0,45 ) e 8 ( 0,45 ), com respectivamente 0,54, 0,49, 0,41 e 0,49 de Índice de Facilidade.
A questão de número 1 apresentou um Índice de Facilidade de 0,30 ( Difícil ), ou seja, menos do que um terço dos estudantes conseguiu resolvê -la, dentro do universo de participantes.
Seu Índice de Discriminação também foi Muito bom ( 0,48 ). nenhuma questão apresentou menos de 15 \% de acertos, razão pela qual seria classificada como Muito difícil.
O Índice de Facilidade variou de 0,30 a 0,86, e o de Discriminação, de 0,41 a 0,51.
A Tabela 6.9 apresenta a distribuição das questões considerando -se simultaneamente a classificação dos dois índices.
As oito questões com Índice de Discriminação Muito bom figuraram entre quatro níveis de dificuldade Muito fácil, Fácil, Médio e Difícil : uma classificada na categoria Muito Fácil ( a questão 6 ) do Índice de Facilidade, duas classificadas na categoria Fácil ( questões 4 e 5 ), quatro na categoria Médio ( questões 2, 3, 7 e 8 ) e, por fim, uma na categoria Difícil ( questão 1 ).
Em particular, a questão 4 foi a que apresentou o maior poder discriminatório, com índice 0,51, tendo sido considerada Fácil em termos de facilidade, com uma proporção de 0,65 acertos.
O máximo de acertos foi alcançado pela questão 6, com um Índice de Facilidade de 0,86.
As questões 4 e 5 foram consideradas Fácil, com índices de facilidade 0,65 e 0,68, respectivamente.
Já os índices de discriminação foram Muito bom para as questões 2 ( 0,41 ), 3 ( 0,45 ), 7 ( 0,45 ) e 8 ( 0,45 ), com respectivamente 0,54, 0,49, 0,41 e 0,49 de Índice de Facilidade.
A questão de número 1 apresentou um Índice de Facilidade de 0,30 ( Difícil ), ou seja, menos do que um terço dos estudantes conseguiu resolvê -la, dentro do universo de participantes.
Seu Índice de Discriminação também foi Muito bom ( 0,48 ).
O Gráfico 6.4, para exemplificar, analisa o comportamento da questão de número 4 de Formação Geral.
Trata -se de uma questão considerada Fácil, com relação à facilidade e a que obteve o maior Índice de Discriminação dessa parte da prova.
Neste gráfico, cada uma das cinco curvas representa o percentual de respostas em determinada alternativa da questão, em função do número de acertos dos estudantes nessa parte da prova ( Formação Geral/Múltipla Escolha ), antes de possíveis eliminações pelo critério do ponto-bisserial.
A curva em verde corresponde à alternativa B, a alternativa correta para essa questão.
Observa -se que entre os estudantes com menor número de acertos, nessa parte do exame, a situação mais frequente foi a escolha da alternativa E ( em vermelho ), incorreta.
Por exemplo, entre os estudantes que acertaram três questões, 25,0 \% escolheram a alternativa E, 43,9 \% escolheram a alternativa correta, B ( em verde ), 18,5 \% escolheram a alternativa D ( roxo ), 7,6 \% escolheram a alternativa A ( azul ) e 4,6 \% a C ( preto ).
Entre os que acertaram três respostas entre as questões de múltipla escolha de Formação Geral, 0,1 \% deixou a questão em branco e 0,2 \% marcou mais de uma alternativa, invalidando a questão.
À medida em que o número de acertos aumenta, indicando desempenho melhor nessa parte da prova, aumenta concomitantemente a proporção de estudantes que selecionaram a alternativa correta B, atingindo 100 \% para os estudantes com oito acertos.
Essa análise permite verificar como a questão discriminou os grupos de desempenho, justificando o alto índice obtido na questão.
Cumpre notar que não é possível inferir deste gráfico nem o índice de facilidade ( que seria uma média da proporção ponderada pela quantidade de alunos com cada uma das notas ), nem o índice de discriminação ponto-bisserial, por razão equivalente.
No caso extremo no qual a grande concentração dos acertos dos alunos fosse abaixo de 4, o índice de facilidade seria obrigatoriamente abaixo de 60 \% ( nesse exemplo ).
Caso a concentração fosse em seis acertos ou mais, o índice seria obrigatoriamente acima de 80 \%.
Os gráficos relativos às demais questões de Formação Geral constam do Anexo I. 6.2.2 Componente de Conhecimento Específico A Tabela 6.10 apresenta as Estatísticas Básicas em relação às questões objetivas do Componente de Conhecimento Específico da prova de Engenharia Civil por Grande Região.
Nessa parte da prova nenhuma das questões foi anulada pela Comissão Assessora da Área.
Além disso, como será discutido ainda nessa seção, oito questões objetivas do Componente de Conhecimentos Específicos não foram usadas no cômputo das notas por terem sido descartadas pelo critério do ponto-bisserial.
Assim, as notas foram calculadas com base em 19 das 27 questões objetivas de Conhecimentos Específicos.
A Média do Brasil deste componente foi de 44,3.
A maior Média foi observada na região Sul ( 48,6 ), e a menor, na região Norte ( 41,3 ).
O Desvio padrão de todo o Brasil foi 15,9, sendo o menor Desvio padrão encontrado na região Norte ( 15,1 ), e o maior, na região Nordeste ( 16,3 ).
A Mediana de todo o Brasil e a de quatro regiões foi 42,1.
Na região Sul a Mediana foi maior que a nacional ( 47,4 ).
A nota Máxima da prova foi 100,0, obtida nas questões objetivas do Componente de Conhecimento Específico, por, pelo menos, um aluno das regiões Sudeste e Sul.
Nas demais regiões a nota Máxima foi 94,7.
A nota Mínima foi zero em todas as regiões.
A Tabela 6.11 apresenta os Índices de Facilidade e Discriminação ( ponto-bisserial ) das questões objetivas do Componente de Conhecimento Específico para os estudantes de Engenharia Civil.
Para facilitar a diferenciação das questões, serão usadas as mesmas cores da Tabela 6.8 para as diferentes classificações dos Índices de Facilidade e de Discriminação.
Como já dito, nenhuma das questões foi anulada pela CAA e as análises dizem respeito às 27 questões válidas.
A partir do Índice de Facilidade obtido, pode -se concluir que um pouco mais de metade das questões objetivas da prova foi considerado, pelo menos, Difícil : das 27 questões, 15 foram classificadas como Difícil ou como Muito difícil.
Três questões foram classificadas como Fácil, e outras nove consideradas como Médio.
Nenhuma questão foi classificada como Muito fácil.
Já quanto aos índices de discriminação das questões objetivas do Componente de Conhecimento Específico da prova, tem -se como resultado a seguinte classificação : seis das 27 questões foram consideradas boas, enquanto duas delas tiveram Índice de Discriminação Muito bom.
Assim, para um pouco menos de um terço das questões – oito em 27 – os Índices de Discriminação foram Bom ou Muito bom.
Dentre as demais, 11 delas foram classificadas como Médio, e outras oito como Fraco, sendo 19, por conseguinte, a quantidade de questões nos dois patamares mais baixos de discriminação.
Constata -se, assim, que a prova – no que se refere ao Componente de Conhecimento Específico, possuía média capacidade de discriminar dentre aqueles que dominavam ou não o conteúdo.
A Tabela 6.12 apresenta a distribuição das questões considerando -se simultaneamente a classificação dos dois índices.
Dentre as questões que alcançaram os maiores índices de discriminação, classificadas com Índice de Discriminação Muito bom, uma delas, a de número 27, foi a de índice mais elevado ( 0,46 ), tendo sido classificada como Difícil quanto ao Índice de Facilidade, com 29 \% dos estudantes marcando a opção correta.
A questão de número 12 foi uma das mais difíceis dentre as 27 questões específicas, com baixo Índice de Facilidade, apenas 12 \% de acertos.
Essa questão apresentou poder discriminatório muito baixo, 0,14, o que comprova ter sido uma das mais difíceis para os estudantes.
Destaca -se, ainda também, a questão 19, também com Índice de Facilidade 0,12, o que, em termos percentuais, corresponde a 12 \% de estudantes que responderam acertadamente.
Já, 0,09 foi o seu Índice de Discriminação.
Tais questões foram, portanto, pelo critério ponto-bisserial, consideradas inadequadas.
Por isso, as questões 11, 17, 23, 25, 30 e 34, além das questões 12 e 19 foram eliminadas do cômputo da nota final.
A Tabela 6.12 apresenta a distribuição das questões considerando -se simultaneamente a classificação dos dois índices.
Dentre as questões que alcançaram os maiores índices de discriminação, classificadas com Índice de Discriminação Muito bom, uma delas, a de número 27, foi a de índice mais elevado ( 0,46 ), tendo sido classificada como Difícil quanto ao Índice de Facilidade, com 29 \% dos estudantes marcando a opção correta.
A questão de número 12 foi uma das mais difíceis dentre as 27 questões específicas, com baixo Índice de Facilidade, apenas 12 \% de acertos.
Essa questão apresentou poder discriminatório muito baixo, 0,14, o que comprova ter sido uma das mais difíceis para os estudantes.
Destaca -se, ainda também, a questão 19, também com Índice de Facilidade 0,12, o que, em termos percentuais, corresponde a 12 \% de estudantes que responderam acertadamente.
Já, 0,09 foi o seu Índice de Discriminação.
Tais questões foram, portanto, pelo critério ponto-bisserial, consideradas inadequadas.
Por isso, as questões 11, 17, 23, 25, 30 e 34, além das questões 12 e 19 foram eliminadas do cômputo da nota final.
A título de exemplo das análises do comportamento das questões objetivas, o Gráfico 6.5 analisa a questão 27 do Componente de Conhecimento Específico.
Esta questão foi considerada pelas respostas dos estudantes avaliados na prova como uma questão com nível de facilidade Difícil, 0,29, ou seja, 29 \% dos estudantes assinalaram acertadamente a opção E, correspondente ao gabarito.
Seu Índice de Discriminação foi igual a 0,46, classificado como Muito bom, o maior valor de discriminação.
Neste gráfico, cada uma das cinco curvas representa o percentual de respostas em determinada alternativa da questão 27, em função do número de acertos dos estudantes nessa parte da prova, antes de possíveis eliminações de questões pelo critério do ponto-bisserial.
A alternativa correta E, representada no gráfico pela curva em vermelho, foi escolhida em maiores proporções pelos alunos com desempenho melhor nessa parte da prova.
Já as alternativas incorretas, também denominadas distratores, foram selecionadas, principalmente, por aqueles com notas mais baixas.
Observa -se que a soma não é 100 \% por causa das questões não respondidas ou com mais de uma opção marcada.
Aqueles com nota zero, na sua quase totalidade, deixaram essa questão em branco ou marcaram mais de uma alternativa, comportamento considerado inválido.
A proporção de alunos que selecionou a resposta correta E aumenta gradativamente, chegando a atingir 100 \% para 23 acertos, enquanto a proporção dos que escolheram alternativas incorretas decai, a partir de três ou quatro acertos, como função do número de acertos nessa parte da prova.
Os gráficos relativos às demais questões do Conhecimento Específico constam do Anexo I. 6.3 ANÁLISE DAS QUESTÕES DISCURSIVAS Esta seção apresenta estatísticas selecionadas e histogramas das Questões Discursivas de Formação Geral ( 6.3.1 ) e Conhecimento Específico ( 6.3.2 ).
São também apresentadas e comparadas as médias de subpopulações caracterizadas por Grande Região. 6.3.1 Componente de Formação Geral As análises dos resultados de desempenho dos estudantes de Engenharia Civil nas duas questões discursivas relativas à Formação Geral encontram -se na Tabela 6.13 e no Gráfico 6.6.
Na Tabela 6.13, observa -se que a nota Média nesse conjunto de questões foi acima da obtida nas objetivas.
Os estudantes, de todo o Brasil, obtiveram, em Formação Geral, Média 55,2 nas questões objetivas e 54,0 nas questões discursivas.
Pode -se notar também que o Desvio padrão nesse conjunto de questões foi um pouco maior do obtido nas objetivas, 21,6 nas questões objetivas e 23,8 nas questões discursivas.
A menor Média foi obtida na região Norte ( 52,4 ), e a maior, na região Sul ( 55,1 ).
A Mediana de todo o Brasil, neste componente, foi 59,5, a mesma obtida nas regiões Nordeste e Sudeste.
Na região Sul a Mediana foi maior ( 60,5 ) e nas demais regiões foi menor que o valor para o Brasil : Norte e Centro-Oeste ( 58,5, em ambas ).
A nota Máxima ( 99,0 ) foi obtida por, pelo menos, um aluno na região Sudeste, sendo : 97,0, nas regiões Norte e Sul, 98,0, na região Nordeste e 97,5, na região Centro-Oeste.
A nota Mínima ( 0,0 ) foi a mesma em todas as regiões do Brasil.
O Gráfico 6.6 representa a distribuição das notas nas questões discursivas do Componente de Formação Geral.
A moda dessa distribuição ocorre no intervalo ( 60 ; 70 ].
Destaca -se também o intervalo [ 0 ; 10 ] com distribuição próxima a 10 \% do total de notas, sendo que no intervalo [ 0 ; 10 ] incluem -se, além da nota zero, a frequência de alunos que deixaram esse tipo de questão em branco.
Na sequência, os resultados verificados para cada uma das questões discursivas de Formação Geral serão apresentados, estabelecendo -se relações com os conteúdos abordados em cada uma delas.
Os comentários da Banca de docentes corretores, a respeito do observado na correção das respostas dos estudantes, suas impressões e conclusões serão apresentados junto à análise de cada questão.
Cumpre esclarecer que, tendo em vista que as questões discursivas de Formação Geral são padronizadas, ou seja, constam de todas as provas, os comentários da Banca são os mesmos para todas as carreiras acadêmicas, sendo direcionados a todos os estudantes que participaram do Enade/2017.
A seguir, serão analisados os desempenhos da Área de Engenharia Civil nas duas questões discursivas de Formação Geral do Enade/2017, comparando -se os resultados obtidos com comentários para cada questão. 6.3.1.1 Análise de Conteúdo da Questão Discursiva 1 do Componente de Formação Geral Os dados de Engenharia Civil, obtidos a partir das respostas à questão 1, encontram -se na Tabela 6.14 e no Gráfico 6.7.
Nessa questão – de melhor desempenho entre as duas de Formação Geral – os alunos, de todo o Brasil obtiveram Média 54,9.
A maior Média para a questão 1 foi obtida na região Sul ( 55,5 ), e a menor, na região Centro-Oeste ( 53,7 ).
Quanto à variabilidade das notas, o Desvio padrão de todo o Brasil foi 27,4.
O menor Desvio padrão foi obtido na região Sul ( 27,0 ), e o maior, na região Centro-Oeste ( 28,1 ).
As medianas de todas as regiões e do Brasil, como um todo, foram iguais ( 60,0 ).
As notas máxima e mínima da questão discursiva 1 foram, sem exceção, as mesmas para todas as regiões do Brasil, respectivamente, 100,0 e 0,0.
O Gráfico 6.7 mostra a distribuição das notas na questão discursiva 1 do Componente de Formação Geral.
Observa -se que a moda foi o intervalo ( 70 ; 80 ], com aproximadamente 22 \% do total de respondentes.
Os alunos que deixaram a resposta a essa questão em branco, em torno de 10 \% dos participantes, representam um máximo local na distribuição das notas. 6.3.1.2 Comentários sobre a correção de Conteúdo das respostas à Questão Discursiva 1 O enunciado era claro e indicava um tema de extrema relevância, qual seja, a incidência de doença grave integrante do rol das DST ( Doenças Sexualmente Transmissíveis ) que encaminha para a análise do direito fundamental à saúde, que tem assento constitucional ( CRFB, art. 6º. caput : São direitos sociais a educação, a saúde, a alimentação, o trabalho, a moradia, o transporte, o lazer, a segurança, a previdência social, a proteção à maternidade e à infância, a assistência aos desamparados, na forma desta Constituição ).
Por um lado, a questão focalizava, por meio de três textos, um tema social extremamente relevante, qual seja, o da relação entre homem e mulher nas suas interações íntimas, no caso representadas pela vida sexual, e o reflexo das relações de poder veiculadas pela sociedade, que apontam, no mais das vezes, o predomínio do homem.
Tal posição de superioridade é justificada pelas convenções arraigadas que acarretam a educação dos homens e mulheres por padrões patriarcais, que redundam em posições machistas.
Por outro lado, um dos textos demandava uma reflexão sobre a realidade dos casais que convivem em relações estáveis ou fugazes e não têm o controle dos efeitos das relações sexuais que abrangem a gravidez não planejada e, em decorrência do não cuidado, a inserção de doenças transmitidas sexualmente.
Esse panorama é refletido na pouca influência das mulheres sobre os homens quanto à proteção, normalmente representada pelo uso de preservativos nas relações intimas, e após a constatação da doença, o encaminhamento as pessoas para tratamento.
Como efeito colateral está a transmissão de doenças para os fetos.
Outro texto trazia à tona o maior cuidado que as mulheres têm, em regra, com sua saúde, notadamente no momento em que se descobrem grávidas, e que redunda na realização de exames obrigatórios de pré-natal que podem facilmente detectar doenças e, caso encontradas, ter encaminhado seu tratamento.
A par disso, destaca -se que a crise de recursos financeiros agravou a situação que já recebe influxos sociais negativos, como a dificuldade do homem em utilizar preservativos, realizar exames e procurar tratamentos quando adoece.
Como muito divulgado pelas mídias, a crise financeira do país repercutiu diretamente na prestação de serviços de saúde em todas as esferas : distrital, federal, estadual e municipal.
Assim, agravou -se o problema da falta de medicamentos essenciais para o tratamento da sífilis e, certamente, de outras doenças.
Como se verifica, a temática indicada no enunciado, representada pelos textos, pertence à realidade conhecida amplamente pelos estudantes, pois também é divulgada por todos os meios de comunicação.
Trata -se de tema geral que, no entanto, normalmente, não se inclui nas grades curriculares de cursos superiores que não sejam da área de Saúde.
Apesar disso, a temática é de conhecimento geral e relacionada aos direitos humanos, estes sim ponto de integração com as diretrizes curriculares.
A linguagem utilizada era clara, indicando a existência dos fatos vinculados a questões relevantes e à saúde pública.
No entanto, como os textos trazem uma ampla gama de possibilidades de interpretação, ocorreram algumas compreensões diversas das esperadas.
Além disso, ressalta -se que a amplitude dos textos disponibilizados no enunciado fornece elementos que poderiam ser simplesmente reproduzidos pelos concluintes nas suas respostas.
O enunciado da questão encaminha a resposta para dois aspectos.
O estudante deveria abordar conteúdo pertinente à vulnerabilidade das mulheres às DST e o papel social do homem em relação à prevenção dessas doenças.
Para responder a esses aspectos, um número expressivo de concluintes utilizou somente os elementos dos textos motivadores, sem nada acrescentar.
A resposta deveria conter também uma segunda parte, a indicação de duas A resposta deveria conter também uma segunda parte, a indicação de duas ações voltadas para o público masculino, a serem adotadas no âmbito das políticas públicas de saúde ou de educação, para reduzir o problema.
Verificou -se que o desempenho dos estudantes nessa questão foi de médio para bom.
Ao analisar as respostas, constata -se que muitos foram os que apresentaram os tópicos presentes no padrão de resposta, amplo o suficiente para albergar as respostas apresentadas pelos concluintes.
Tal resultado pode ser justificado, tendo em vista tratar -se de tema bastante difundido nas diversas mídias, apresentado de forma ampla pelos textos motivadores presentes no enunciado.
Por outro lado, o debate no ambiente universitário sobre as relações de poder entre homens e mulheres tende a ser marcante pelo movimento de empoderamento feminino, liderado por organizações sociais com ampla atuação nos estabelecimentos de nível superior.
De forma recorrente, as relações patriarcais entre homens e mulheres foram vinculadas ao modelo machista vigente.
Essas condições, segundo os concluintes, explicam a incidência da doença e a situação de vulnerabilidade das mulheres, bem como a função social dos homens.
Outras respostas apontaram o poder público, nas áreas de saúde e/ou educação, como órgãos responsáveis por resolver o problema.
Para os que se fixaram na área da saúde, foram frequentes a indicação da necessidade de mais medicamentos, vacinas, educação sexual, médicos e postos de atendimentos.
Alguns, no entanto, apenas indicaram ações individuais para o homem, e outros, ações individuais para as mulheres.
Também figuraram respostas, apresentando como solução, ações da igreja e fazendo referências a campanhas no seio das famílias. 6.3.1.3 Análise de Conteúdo da Questão Discursiva 2 do Componente de Formação Geral A Tabela 6.15 mostra que o desempenho médio dos estudantes na questão discursiva 2 ( média 52,4 ) foi inferior ao obtido na questão discursiva 1 ( média 54,9 ).
A região Sul foi aquela cuja Média, também nessa questão, foi maior ( 53,9 ), enquanto a de menor Média foi a região Norte ( 49,8 ).
Quanto à variabilidade das notas, o Desvio padrão de todo o Brasil foi 32,4, superior ao obtido na questão discursiva 1 ( 27,4 ).
O maior desvio nessa questão foi obtido na região Norte ( 33,3 ), enquanto o menor foi obtido na região Sudeste ( 32,0 ).
A Mediana de todo o Brasil foi 50,0, a mesma de todas as regiões, exceto na região Sul ( 70,0 ).
As notas máximas ( 100,0 ) e as notas mínimas ( 0,0 ) também foram as mesmas em todas as regiões, sem exceção.
O Gráfico 6.8 mostra a distribuição das notas na questão discursiva 2 do Componente de Formação Geral.
Observa -se que a maior frequência corresponde aos alunos que obtiveram nota dentro do intervalo ( 70 ; 80 ].
Destacam -se, também, como um máximo local, os alunos que deixaram a questão em branco.
Se estes forem somados àqueles que, tendo respondido a questão, obtiveram nota zero, chega -se a cerca de 20 \% dos respondentes.
Outro máximo local é o intervalo ( 90 ; 100 ], com mais de 10 \% dos estudantes.
Nota -se, ainda, que as notas ficaram mais dispersas em comparação à questão discursiva de número 1.
Isso pode ser constatado também pela comparação do desvio padrão das notas da questão discursiva 2 ( 32,4 ) e da questão discursiva 1 ( 27,4 ). 6.3.1.4 Comentários sobre a correção de Conteúdo das respostas à Questão Discursiva 2 O enunciado era claro e abordava um tema de extrema relevância : a questão de gênero.
Trata -se de uma proposta polêmica, pois abrange a análise da inserção social de minoria vinculada à questão de gênero.
O enunciado solicitava que o concluinte discorresse sobre a importância do nome para as pessoas transgêneras e propusesse uma medida, no âmbito das políticas públicas, para facilitar o acesso dessas pessoas à cidadania.
A respeito da primeira parte, a importância do nome, esperava -se que o estudante abordasse as dificuldades que o uso do nome civil ( nome de registro, nome de batismo, nome da certidão de nascimento, por exemplo ) pode acarretar para as pessoas transgêneras, além de demonstrar conhecimento sobre a importância do nome social, direito já reconhecido através, principalmente, de decisões judiciais e atos administrativos, como decretos e portarias que dispõem sobre o uso do nome social.
Em atendimento à segunda parte do comando, o participante deveria propor uma política pública associada ao direito à cidadania das pessoas transgêneras.
Sobre esse aspecto, esperava -se a compreensão do significado de política pública, ou seja, propostas de âmbito governamental e não iniciativas de Organizações Sociais, igrejas ou grupos sociais de qualquer tipo.
O nascimento das pessoas, no Brasil, é objeto de inscrição em cartórios de Registro Civil das Pessoas Naturais, órgão responsável por estabelecer a certificação pública dos dados de todos os indivíduos que não somente nascem, mas casam, eventualmente se divorciam e vêm a óbito.
Daí o nome civil.
Para flexibilizar a correção, outras nomenclaturas foram consideradas, a saber : identificação de nascimento ; nome de batismo ; nome de registro, nome próprio, nome verdadeiro, que diferencia o nome civil daquele que o indivíduo deseja utilizar para traduzir a sua opção pessoal de gênero, que se convencionou denominar nome social.
Quanto ao entendimento do termo transgênero, deve -se destacar que o movimento de mudança do gênero é contemporâneo e não está necessariamente vinculado à mudança de sexo.
Existem várias hipóteses em jogo.
Uma é a que envolve cirurgias de mudança de sexo ( a menina extirpa os seus seios e inclui órgão genital por intervenção cirúrgica ; o menino inclui seios e extirpa o seu órgão genital e inclui o órgão feminino também por intervenção cirúrgica ).
Há outro movimento que envolve apenas as vestimentas e o linguajar corporal, além dos relacionamentos ( as meninas assumem o trajar, o andar, os gostos dos meninos e mantêm relacionamentos afetivos com meninas ; os meninos assumem o trajar, o andar, os gostos das meninas e mantêm relacionamentos afetivos com meninos ).
Em relação à discussão da primeira solicitação do comando da questão, destaca -se que todos têm nome civil, que foi escolhido pelos seus pais ou responsáveis e, na falta deles, por autorização do Juiz de Direito, pautado no sexo detectado no nascimento.
Já o nome social é aquele que indica a preferência do indivíduo na sua identificação nas relações sociais e que destoa do seu gênero de registro civil.
Ressalte -se que o Brasil não possui ainda legislação específica sobre pessoas transgêneras.
Observa -se um relativo progresso no reconhecimento de direitos, através, principalmente, de decisões judiciais e atos administrativos, como decretos e portarias que dispõem sobre o uso do nome social.
Sem uma lei que defina os procedimentos da alteração dos documentos para pessoas transexuais, essa parcela da população LGBT ( lésbicas, gays, bissexuais, transexuais e transgêneros ) é obrigada a procurar na Justiça o reconhecimento de sua identidade, em processos que podem ser longos e que dependem do olhar de determinados profissionais como médicos, psicólogos, assistentes sociais, bem como do entendimento dos advogados, defensores públicos, juízes e do Judiciário como um todo.
Atualmente, tramita na Comissão de Direitos Humanos da Câmara dos Deputados o Projeto de Lei João Nery ( 5002/2013 ), dos deputados Jean Willys ( PSOL-RJ ) e Erika Kokay ( PT-DF ), que determina que o reconhecimento da identidade de gênero seja um direito do cidadão.
O projeto recebeu o nome do primeiro trans-homem operado no Brasil.
Grande parte dos respondentes considerou a incongruência entre a imagem corporal, representada pelo gênero assumido pelos trans, e o nome que apresenta nos seus documentos como causa de sérios constrangimentos e embaraços e as dificuldades para que esses sujeitos acessem direitos básicos como saúde, educação e empregabilidade em uma sociedade heteronormativa e preconceituosa.
Muitos demonstraram reconhecer que a identidade e a expressão de gênero, através do nome social, pressupõe que sejam respeitadas a dignidade da pessoa humana, a intimidade, a liberdade e o direito à cidadania plena.
No entanto, também foram encontradas, em quantidade significativas, respostas que demonstravam a existência de uma interpretação social conservadora que não admite a mudança de sexo e nem do nome social.
Existe um acentuado preconceito contra todos os movimentos que buscam identificação social diversa da civil.
Isso inclui os variados grupos : gays, travestis, transgêneros, dentre outros.
Já em relação ao segundo tópico a ser abordado pelo concluinte, é notório que as pessoas trans estão destituídas de direitos no campo das políticas públicas como saúde, habitação, emprego, educação, justiça, entre outros.
Nesse sentido, uma grande diversidade de medidas, no âmbito das políticas públicas, poderia ser proposta.
Ações voltadas para a capacitação de profissionais da educação, saúde, segurança pública entre outras áreas, para melhor atender às demandas e especificidades desse público alvo seriam necessárias.
Ainda em relação à educação, também seriam necessárias ações que contribuíssem para diminuir a evasão, pois o ambiente escolar é difícil para uma pessoa transgênera, por conta das situações discriminatórias e preconceituosas vindas dos colegas de classe, e até mesmo das/dos professores e/ou de outros profissionais da área, que não podem ser mais um complicador nesse processo.
Ainda no que tange ao assunto, percebe -se que ao serem excluídas dos bancos escolares, as oportunidades no mercado de trabalho formal praticamente inexistem.
Um outro aspecto que se evidencia no cotidiano desse segmento da sociedade diz respeito às situações de violência sofridas pelas pessoas transgêneras.
Nesse sentido, algumas propostas foram elencadas para minimizar o problema e combater a violência.
Políticas envolvendo uma melhor formação de funcionários das polícias civil e militar, por exemplo, somariam esforços na qualificação das abordagens policiais e dos atendimentos prestados pelas forças de segurança, pois o respeito à orientação sexual e à identidade e expressão de gênero é antes de tudo o respeito à dignidade, à cidadania e à própria democracia.
Nesse sentido, revela -se essencial que, ao sofrerem agressões, sejam elas físicas, psicológicas ou morais no espaço público ou domiciliar, as pessoas trans possam ser atendidas, com respeito inclusive ao nome social, nas delegacias, inclusive nas Delegacias Especiais de Atendimento à Mulher ( DEAM’s ), pois esse é o órgão institucional que possivelmente detém as qualificações necessárias para o atendimento ao gênero feminino ao qual as travestis e mulheres trans efetivamente pertencem.
Também ocorre uma latente inacessibilidade das pessoas trans à assistência jurídica.
Assim, é necessário implementar projetos de capacitação para profissionais do poder judiciário como um todo, de modo a prestarem a assistência jurídica gratuita na defesa intransigente dos direitos sociais, civis, trabalhistas e políticos dessa população.
Dessa maneira, a Justiça agiria de forma mais acessível e sensível para o segmento, que terá suas especificidades observadas a começar pelo devido respeito do direito ao uso do nome social.
Nesse âmbito, o mais comum foi propor a adoção de medidas punitivas para quem viola o direito à autodeterminação de gênero e à elaboração de leis que garantam a mudança do nome e assegurem outros direitos para as pessoas trans.
Propostas envolvendo a formação de profissionais, como as mencionadas aqui, e de conscientização da população foram as mais frequentes.
Muitas outras enfatizavam o direito ao uso respeitado do nome social e à garantia de alteração do registro civil, independentemente da cirurgia de redesignação sexual ou transgenitalização.
Muitos questionaram o fato de não bastar a autodeclaração no processo de alteração de nome, enfatizando que o nome social é uma questão de princípio de dignidade e cidadania.
Apesar de pouco mencionada nas respostas corrigidas, políticas envolvendo a promoção da saúde são muito necessárias.
Cabe aos gestores de saúde ações direcionadas, não somente à prevenção e ao tratamento do HIV/Aids, mas também políticas de saúde que contemplem as especificidades dessa população, tais como o Ministério da Saúde, especialmente o SUS ( Sistema Único de Saúde ), incluir, de fato, o atendimento para pessoas trans para acompanhamento da terapia hormonal que muitas demandam, respeitando o direito de utilização do nome social.
De modo geral, são muitas as políticas públicas que precisam ser implemantadas e urge a necessidade de uma agenda mais compromissada por parte dos gestores do nosso país com essa comunidade, que ainda é muito discriminada.
Logo, tais políticas devem contemplar suas necessidades básicas, como o direito de acesso aos estudos, à capacitação para o mercado de trabalho, à profissionalização, ao acesso a bens e serviços de qualidade em saúde, habitação, segurança, cultura e assistência social e, certamente, o respeito ao uso do nome social, o que também contribui para essa inclusão.
Assim, as opções apresentadas no padrão de resposta funcionaram como exemplos de políticas públicas possíveis, já que seria difícil apresentar uma lista suficientemente ampla para abranger a grande variedade de respostas encontradas.
O desempenho dos estudantes nessa questão foi médio.
As respostas, de modo geral, indicaram a necessidade de uma atualização por parte dos concluintes quanto a temas essenciais e, no caso da presente questão, a análise sobre aspectos fundamentais da questão de gênero.
Muitos, possivelmente motivados pelos textos apresentados no enunciado, restringiram -se a discorrer sobre preconceito de gênero ou mesmo sobre o preconceito de modo mais geral, sem responder ao que foi solicitado no comando da questão.
Outros apresentaram discursos religiosos ou ideológicos, também sem tratar do que foi solicitado. 6.3.1.5 Análise de Língua Portuguesa das Questões Discursivas do Componente de Formação Geral Os dados de Engenharia Civil, obtidos a partir das respostas às questões discursivas do Componente de Formação Geral, no que tange à Língua Portuguesa, encontram -se na Tabela 6.16 e no Gráfico 6.9.
Nesse aspecto, os alunos, de todo o Brasil, obtiveram Média 55,5.
A maior Média com respeito à Língua Portuguesa foi obtida na região Sul ( 56,6 ), e a menor, na região Norte ( 53,8 ).
Quanto à variabilidade das notas, o Desvio padrão de todo o Brasil foi 21,7.
O menor Desvio padrão foi obtido na região Sudeste ( 21,2 ) e o maior Desvio padrão foi obtido na região Centro-Oeste ( 22,7 ).
A Mediana das notas de Língua Portuguesa foi 62,5 para o Brasil como um todo, a mesma obtida em três regiões.
As exceções foram as Medianas das regiões Norte e Centro-Oeste, iguais a 60,0.
A nota Máxima para todo o Brasil foi de 97,5, com, pelo menos, um aluno tirando essa nota na região Sudeste.
Nas demais regiões a nota Máxima foi 95,0.
Já a nota Mínima foi zero em todas as regiões do país.
O Gráfico 6.9 mostra a distribuição das notas de Língua Portuguesa do Componente de Formação Geral.
Observa -se que a maior frequência ( quase 30 \% ) corresponde aos alunos que obtiveram nota no intervalo ( 60 ; 70 ].
Destaca -se, também, os alunos que deixaram ambas as questões em branco representando um pouco menos do que 10 \% do total, caracterizando -se como máximo local. 6.3.1.6 Comentários sobre a correção das respostas de Formação Geral com respeito à Língua Portuguesa Os enunciados apresentados em 2017 não explicitaram a exigência de elaboração de um texto “ dissertativo ”, como nos anos anteriores.
O participante deveria fazer uma exposição de seus conhecimentos sobre cada um dos assuntos e estruturar seus textos de acordo com as características do registro formal adequado à situação comunicativa – avaliação de conhecimentos.
Essa configuração determina exigências quanto aos seguintes aspectos : adequação da seleção vocabular, desenvolvimento do conteúdo, organização lógica das ideias, estruturação sintática dos períodos, utilização de procedimentos de encadeamento textual e de referenciação, obediência às exigências morfossintáticas próprias da modalidade escrita da norma-padrão, respeito às regras ortográficas e às regras de acentuação gráfica.
O padrão de resposta utilizado na avaliação das questões 1 e 2 considerou os aspectos relevantes ao bom desempenho linguístico como competências distintas, de modo a permitir um mapeamento detalhado do domínio dos recursos disponíveis na Língua Portuguesa para a comunicação escrita formal.
Com base nesse objetivo, foram avaliados os seguintes aspectos : ( a ) estruturação textual condizente com o gênero solicitado e o modo de organização textual expositivo adequado ao gênero – essa competência envolve : • estruturação sintática condizente com o padrão da modalidade escrita formal da língua portuguesa de modo a garantir a clareza necessária ; • distribuição do conteúdo do texto em parágrafos, de modo a garantir a sua organização temática ; • utilização de operadores discursivos que contribuam para a progressão temática do texto, estabelecendo relações lógicas entre as ideias apresentadas, tanto do ponto de vista intrafrasal, como do interfrasal ; • utilização de procedimentos de referenciação lexical e pronominal que permitam a retomada de referentes textuais ; • utilização de sinais de pontuação que contribuam para a organização lógica da frase e do texto ; • inteligibilidade relacionada ao atendimento das exigências de estruturação textual.
Espera -se, portanto, que o participante recorra a procedimentos linguístico-discursivos para organizar seu texto, permitindo o encadeamento lógico entre suas partes de forma a garantir a progressão e a coerência textuais.
Isso significa que os seguintes procedimentos foram considerados inadequados, de acordo com o padrão de resposta proposto :  elaboração de frases fragmentadas que comprometam a estrutura lógico-gramatical do texto ;  sequência justaposta de ideias sem encaixamentos sintáticos, reproduzindo hábitos da oralidade ;  elaboração de frase com apenas oração subordinada, sem oração principal ;  emprego equivocado do conector ( preposição, conjunção, pronome relativo, alguns advérbios e locuções adverbiais ) comprometendo a expressão da relação lógica entre duas ideias, com prejuízo da clareza do texto ;  emprego do pronome relativo sem a preposição, quando obrigatória ;  repetição ou substituição inadequada de palavras sem o emprego dos recursos oferecidos pela língua ( pronome, advérbio, artigo, sinônimo ) ;  emprego inadequado dos pronomes relativos “ cujo(a) ” e “ onde ” ;  utilização inadequada dos sinais de pontuação, comprometendo a clareza textual ;  inteligibilidade relacionada ao atendimento das exigências de estruturação textual. ( b ) respeito às convenções ortográficas da norma-padrão da Língua Portuguesa – essa competência envolve o domínio das regras de acentuação gráfica e da grafia padrão das palavras ( com ausência de abreviaturas próprias da linguagem da internet ), de acordo com as convenções estabelecidas pela legislação em vigor e consubstanciadas no Vocabulário Ortográfico da Língua Portuguesa, editado pela Academia Brasileira de Letras. ( Vale notar que, nessa edição de 2017, não foi aceita a legislação anterior, no caso das regras relativas ao uso do hífen e da acentuação gráfica ).
Espera -se que o participante :  grafe corretamente as palavras ;  respeite as regras de acentuação gráfica ;  empregue maiúsculas em início de frase, em nomes próprios de pessoas, lugares ou instituições ;  grafe as siglas com letras maiúsculas ou apenas com a primeira letra maiúscula quando formarem uma palavra ;  evite abreviações como p/, vc, tb, pq, tá, né, usadas muitas vezes em escrita informal e na internet ;  obedeça às regras de separação de sílabas no final da linha. ( c ) domínio dos diferentes aspectos morfossintáticos próprios da modalidade escrita formal da norma-padrão da Língua Portuguesa – essa competência envolve : a concordância nominal, a concordância verbal, a regência nominal, a regência verbal, a flexão nominal, a flexão verbal, a correlação entre os tempos verbais, a colocação pronominal e a utilização de sinais de pontuação que contribuam para a organização lógica da frase e do texto.
Espera -se que o participante :  flexione o verbo para estabelecer concordância de número com o sujeito da frase ;  flexione o artigo, o adjetivo e o pronome para concordar em número e em gênero com o substantivo a que se referem ;  observe a regência nominal e a verbal, utilizando a preposição adequada depois de um substantivo, um verbo ou um adjetivo ;  empregue adequadamente o acento grave indicador da crase entre uma preposição e um artigo ( a+a ) ;  obedeça às regras de colocação pronominal ( próclise e ênclise ), distintas dos hábitos da oralidade ou da escrita informal ( exigência de próclise com termo atrator, não exigência de que o pronome oblíquo se ligue ao verbo auxiliar por meio de hífen ) ;  flexione adequadamente verbos, substantivos, adjetivos e pronomes no que diz respeito à expressão das categorias gramaticais ;  flexione os verbos para expressar a correlação de modo e tempo nas estruturas subordinadas.
Com base nesses critérios, foram considerados como desvios de caráter morfossintático, e não como desvios ortográficos, as alterações que envolvem mudança de classe gramatical ou de forma flexional do verbo :  eliminação da marca de infinitivo ( -r- ) e substituição por acento agudo ou ausência total de marca do infinitivo ;  confusão entre “ ão ” e “ am ” nas formas verbais ;  confusão entre “ há ” e “ a ” ;  uso de hífen para separar pronome átono – tanto uso indevido quanto omissão ( exemplo : “ esperasse ”, em lugar de “ espera -se ” ; “ falar -mos ”, no lugar de “ falarmos ” ;  “ esta ” ( no lugar de “ está ” ) ; “ mais ” ( no lugar de “ mas ” ) ; “ e ” ( no lugar de “ é ” ).  verbos “ ter ” e “ vir ” que, na terceira pessoa do plural, não apresentarem o acento circunflexo, foram considerados como desvio de concordância, em aspectos morfossintáticos. ( d ) seleção vocabular adequada à modalidade escrita formal da Língua Portuguesa, exigida pela situação comunicativa – essa competência envolve a precisão na seleção/utilização do vocabulário relacionado à temática solicitada pela questão ; a ausência de marcas da oralidade, como termos de sentido muito genérico ( “ coisa ”, “ negócio ”, “ você ” ) e termos de registros mais informais ( como gírias, jargões, frases feitas, ditados populares, termos regionais ).
Assim, espera -se que o participante respeite a adequação vocabular não usando gírias ou expressões coloquiais, evite repetição desnecessária de palavras e utilize um vocabulário mais formal, como solicitado por um texto dissertativo.
Observações : - A inteligibilidade relacionada ao atendimento das exigências de estruturação textual foi avaliada na segunda competência ( relativa aos aspectos textuais ). - Os problemas de coerência textual provocados por uso indevido do vocabulário foram avaliados na terceira competência ( relativa aos aspectos morfossintáticos e vocabulares ). - Não foram considerados como desvios morfossintáticos os problemas de caligrafia ( -a/ -o, -s, -r ). - Os textos grafados integralmente em caixa alta foram corrigidos da mesma forma que os demais, já que não seria possível distinguir alguma marcação especial para as letras em início de frase. - Cada desvio foi considerado como uma ocorrência, mesmo que dois desvios fossem relativos ao mesmo aspecto linguístico.
A grade de avaliação do desempenho linguístico considerou, portanto, três grandes grupos de competências, segundo os aspectos explicitados anteriormente : a ) domínio das convenções ortográficas : grafia de vogais e consoantes, uso de maiúsculas e minúsculas, emprego do hífen e acentuação gráfica ; b ) domínio dos procedimentos de estruturação textual do ponto de vista microestrutural : organização interna dos períodos, emprego de conectores para a articulação lógica entre os períodos e entre os parágrafos, emprego de marcas de referenciação lexical e pronominal ; utilização dos sinais de pontuação que contribuem para a organização lógica da frase ; c ) domínio das regras de caráter morfossintático estabelecidas como modelares do ponto de vista da modalidade escrita formal da norma-padrão da Língua Portuguesa : concordância nominal e verbal, regência nominal e verbal, colocação pronominal, flexão nominal e verbal, correlação entre tempos e modos verbais, ausência de marcas de oralidade.
A seleção vocabular adequada à modalidade escrita formal da Língua Portuguesa foi incorporada a essa última competência, tendo em vista a intersecção entre as duas do ponto de vista das exigências do registro formal da modalidade escrita da norma-padrão.
Apreciação geral do desempenho dos estudantes : A correção revelou desempenhos distintos dos participantes, marcados pela falta de repertório cultural da maioria dos participantes : - a questão 1 propiciou a oportunidade de maior desenvolvimento do tema solicitado, gerando consequências na estruturação textual, já que os textos foram mais longos, mais elaborados e fluentes.
Entretanto, em virtude da diversidade de abordagens dos textos motivadores ( a falta de penicilina, a recusa no uso da camisinha pelos homens, a maior frequência das mulheres na realização de exames ginecológicos, a transmissão das DST aos bebês durante a gestação ), observaram -se respostas que se configuraram como verdadeiras paráfrases, com comprometimento do caráter autoral desejável em uma questão discursiva de Formação Geral.
Alguns participantes selecionaram um trecho de cada texto motivador para compor suas respostas, outros copiaram trechos inteiros. - a questão 2 teve um comportamento distinto, devido ao caráter polêmico que o caracteriza.
Assim, os textos são mais curtos, com vocabulário repetitivo e preso aos textos motivadores.
Além disso, observou -se grande ocorrência de protesto em função do tema, ora por motivos religiosos ora por preconceito de gênero.
Quanto aos aspectos linguísticos analisados durante esta avaliação, os resultados observados estão descritos nos parágrafos seguintes.
Aspectos ortográficos : O desempenho dos participantes revelou uma diferença muito grande nos dois aspectos analisados nesta competência : baixo índice de desvios da grafia padrão e grande índice de desvios de acentuação.
Em vários casos, ocorre ausência completa de acentuação gráfica.
Os resultados revelam que a tendência dominante entre os universitários brasileiros é a eliminação da acentuação gráfica, talvez motivada pelos hábitos relacionados às redes sociais e pela ausência de esclarecimento dos meios de comunicação, das autoridades e das escolas sobre as decisões do Acordo Ortográfico da Língua Portuguesa de 1990.
Os casos mais sistemáticos de eliminação do acento indicador da sílaba tônica são :  palavras proparoxítonas ( “ sifilis ”, “ proximos ”, “ politicas ”, “ publicas ”, “ transgeneros ” ) ; Aspectos textuais : Esta competência é a que se revela como a mais problemática entre os participantes, porque são muitos os problemas observados, desvios acumulados durante toda a formação do estudante e que não se resolvem com um estudo autodidata, como acontece com regras ortográficas ou morfossintáticas.
São eles : sequência justaposta de ideias sem encaixamentos sintáticos ; redução drástica de estruturas subordinadas, ao lado do aumento na frequência de estruturas coordenadas e absolutas ; redução no uso de conectores para expressar relações lógicas essenciais à construção do texto, substituídas pela exigência de inferência por parte do interlocutor para suprir a sua ausência ; emprego equivocado de operadores que não estabelecem relações lógicas coerentes entre ideias do texto ; emprego inadequado do pronome relativo ( com omissão da preposição ou a utilização de pronome inadequado, como “ onde ” ) ; repetição exaustiva de termos sem a utilização de procedimentos mais sofisticados de substituição ( hiperonímias, hiponímias, nominalizações, expressões  palavras paroxítonas terminadas em ditongo crescente ( “ necessario ”, “ noticiarios ”, “ individuo ”, “ dependencia ”, “ varios ”, “ propria ”, “ transmissiveis ”, “ ocorrencia ” ) ;  palavras paroxítonas com hiato ( “ saude ” )  palavras oxítonas ( “ ninguem ”, “ esta ”, “ ate ”, “ tambem ”, “ prevencao ” ).
Por outro lado, destaca -se o uso indevido do acento gráfico em palavras como “ genêro ”, “ melâncolia ”, “ prevênção ”, “ intervênção ”, por exemplo.
Há, também, uma tendência a acentuar paroxítonos como se fossem oxítonos : “ transmissivéis ”.
Quanto ao domínio das convenções relativas à grafia das palavras, observam -se desvios recorrentes, como a hipercorreção pela escolha de “ e ” no lugar de “ i ”, por influência de hábitos da oralidade ( “ descriminação ” no lugar de “ discriminação ”, “ entervenção ” no lugar de “ intervenção ” ) ; desvios de grafia relacionados à variação diastrática podem ser observados em “ estrupo ”, “ subjulgadas ”, “ precoseito ”, “ soubre ”, “ vecendo ”, “ indesencia ”, “ apolojia ”, “ fulga ”, “ dereitos ”, “ sifelis ”.
Observam -se, também, casos de inadequação no uso da maiúscula para destacar determinadas palavras-chave do texto, como “ Transgêneros ”, “ Brasileiros ”, “ Homens ”.
Destaque -se, também, o grande número de participantes que grafam os textos inteiramente em caixa alta.
Vale observar, também, que, ao contrário do que se esperava, não apareceram abreviaturas próprias do “ internetês ”, relacionadas ao uso de redes sociais e e-mails.
Aspectos textuais : Esta competência é a que se revela como a mais problemática entre os participantes, porque são muitos os problemas observados, desvios acumulados durante toda a formação do estudante e que não se resolvem com um estudo autodidata, como acontece com regras ortográficas ou morfossintáticas.
São eles : sequência justaposta de ideias sem encaixamentos sintáticos ; redução drástica de estruturas subordinadas, ao lado do aumento na frequência de estruturas coordenadas e absolutas ; redução no uso de conectores para expressar relações lógicas essenciais à construção do texto, substituídas pela exigência de inferência por parte do interlocutor para suprir a sua ausência ; emprego equivocado de operadores que não estabelecem relações lógicas coerentes entre ideias do texto ; emprego inadequado do pronome relativo ( com omissão da preposição ou a utilização de pronome inadequado, como “ onde ” ) ; repetição exaustiva de termos sem a utilização de procedimentos mais sofisticados de substituição ( hiperonímias, hiponímias, nominalizações, expressões metafóricas ) ; frases fragmentadas que comprometem a estrutura lógico-gramatical ; frases formadas apenas por oração subordinada, sem oração principal.
Um importante aspecto a destacar é o baixo desempenho de uma parte dos participantes em relação à estrutura formal do texto produzido, o que é preocupante ao se levar em conta que são graduandos em fase final de formação.
São frequentes os casos de desvios de estruturação frasal, com uso inadequado ou ausência de conectivos entre parágrafos e entre frases.
Em uma parte dos textos, falta textualidade e domínio do registro padrão da língua.
Na verdade, observam -se relações linguísticas quase agramaticais, como as estabelecidas pela sequência de gerúndios sem o apoio de um ponto de partida para a organização das informações gramaticais e semânticas.
Observou -se que uma parte dos participantes não distribuiu as ideias em parágrafos, talvez devido ao pequeno número de linhas disponibilizadas para a resposta da questão ou, quem sabe, pela suposição de que não seria necessária essa divisão por não se tratar de um texto no modelo de uma redação dissertativo-argumentativa, como solicitado nos vestibulares.
Em função do encaminhamento dos enunciados das duas questões, que solicitaram uma análise do problema e encaminhamentos de políticas públicas, houve uma grande tendência, também, de construção de dois parágrafos desconexos, sem utilização de elementos coesivos adequados à progressão textual.
Quanto à utilização dos mecanismos de referenciação, deve -se destacar a ocorrência de repetições de palavras ou expressões sem a utilização de termos sinônimos ou pronomes, como seria adequado.
Quanto à utilização dos sinais de pontuação, observou -se uma grande precariedade nos textos analisados.
É muito frequente a ocorrência de parágrafos sem marca interna de pontuação para separar os períodos.
Vale observar que não foi penalizada a ausência de vírgula para destacar locuções ou adjuntos adverbiais de pequena extensão deslocados de posição na frase, por ser um uso opcional.
São os seguintes os tipos de problemas encontrados : a ) vírgula : utilização de vírgula para separar o sujeito e o predicado ; ocorrência de apenas uma das vírgulas para separar uma palavra, uma expressão ou uma oração encaixada ; uso de vírgula no lugar do ponto para separar ideias que constituem períodos distintos ; ausência de vírgula para separar elementos de uma enumeração ; ausência de vírgula para separar oração adjetiva explicativa ou utilização inadequada para separar oração adjetiva restritiva ; b ) ponto e vírgula : utilização do ponto e vírgula no lugar de vírgula ; c ) ponto final : ausência de ponto final para separar períodos.
Aspectos morfossintáticos e vocabulares : Em relação à regência, o desvio mais frequente é a falta do sinal indicativo da crase – isso revela que o usuário não tem consciência de que, sob a forma do termo “ a ”, existe a presença de uma contração entre a preposição “ a ” ( exigida pela regência do termo anterior ) e o artigo definido “ a ”.
Um desvio de regência significativo, nos últimos anos, é a utilização inadequada de uma preposição ou sua ausência após o verbo ou o nome ( substantivo ou adjetivo ).
Outro problema relacionado à regência verbal e à nominal, encontrado frequentemente nas questões, foi a ausência de preposição antes de pronome relativo, processo generalizado na modalidade oral da língua, em situações de registro informal.
Apesar da possibilidade de que essa alteração de regência se generalize no padrão escrito da Língua Portuguesa, como já está ocorrendo até em textos jornalísticos, o não emprego da preposição foi considerado inadequado neste processo de avaliação.
A concordância verbal e a concordância nominal apresentam alguns desvios muito frequentes.
Quanto à concordância de número, observou -se ausência de marca ( com sujeito anteposto ou posposto ) ou uso indevido ( uso inadequado da marca de plural comandado pelo núcleo plural da locução adjetiva, apesar de o substantivo que funciona como núcleo do sintagma nominal estar no singular ).
Uma ocorrência que se destacou foi a ausência de acento circunflexo na forma plural do presente do indicativo do verbo “ ter ”, que foi considerada como um desvio na concordância verbal e não na acentuação gráfica.
Quanto à concordância de gênero, vários casos foram observados, normalmente no âmbito de sintagmas nominais longos, em que o adjetivo está afastado do substantivo.
Deve -se destacar o aparecimento da marca de plural em verbos ou adjetivos comandados por adjunto adnominal plural, apesar de serem relacionados a núcleos substantivos no singular, evidenciando um processo de hipercorreção ( exemplo : “ O nome social dos trangêneros podem ser estabelecidos por uma legislação específica ” ).
Quanto à questão da colocação pronominal, foram poucos os desvios observados.
Concluiu -se que, no registro escrito formal, a maioria dos participantes já incorporou regras como a não introdução da frase por um pronome oblíquo e a próclise na presença de um termo atrator.
Não se adotou, entretanto, o padrão excessivamente formal descrito pelas gramáticas normativas em relação à posição do pronome oblíquo em locuções verbais, já que esse uso está muito distante da prática cotidiana, até em textos mais formais.
Quanto aos aspectos vocabulares, alguns tipos de inadequação foram observados : expressões da oralidade ; seleção vocabular incompatível com o contexto, gerando falta de inteligibilidade ; falta de domínio de vocabulário mais abstrato e de maior complexidade, essencial ao desenvolvimento do texto de base dissertativa.
O principal aspecto observado foi a excessiva repetição de certas palavras, revelando limitação de repertório vocabular.
O termo “ pessoa ”, por exemplo, chega a ser repetido até 7 ou 8 vezes em um mesmo texto.
Várias marcas de oralidade foram identificadas, embora não com alta frequência : uso do pronome relativo “ onde ” como relativo universal, falta de artigo definido antes de substantivo, reduções como “ tá ”, “ pra ”, “ pro ”, “ prum ”, expressões informais, eliminação de preposições.
Em função do tema solicitado na questão 1, a sigla DST foi grafada de diferentes maneiras : DST´s, DSTs ou DSTS.
Além disso, houve flutuação de gênero gramatical nos sintagmas que envolviam a sigla : “ os diferentes DSTs ” ou “ as diferentes DSTs ”.
Em função do tema solicitado na questão 2, houve flutuação de gênero gramatical nos sintagmas que envolviam a sua palavra-chave : “ os trans ” ou “ as trans ” ; “ os transgêneros ” ou “ as transgêneras ”.
Todas essas formas foram aceitas, por considerarmos que sua utilização é muito recente e os participantes ainda não sistematizaram essas formas em seu uso da língua. 6.3.2 Componente de Conhecimento Específico Na parte da prova relativa às questões discursivas do Componente de Conhecimento Específico ( Tabela 6.17 ), observa -se que a Média foi mais baixa do que para as questões discursivas do Componente de Formação Geral.
Enquanto no Componente de Formação Geral, a Média para estudantes de Engenharia Civil de todo o Brasil foi 54,0, na parte de Conhecimento Específico, a Média foi 18,8.
A maior Média desse componente foi obtida pelos estudantes da região Sul ( 22,5 ), e a menor, pelos da região Sudeste ( 17,0 ).
Quanto à variabilidade das notas, o Desvio padrão de todo o Brasil foi 15,7.
O maior Desvio padrão foi encontrado na região Nordeste ( 17,0 ), e o menor, nas regiões Norte e Sudeste ( 14,9 ).
A maior nota Máxima, 100,0, foi obtida em todas as regiões do Brasil, exceto na região Norte ( 91,7 ).
A nota Mínima ( 0,0 ) foi obtida por pelo menos um aluno em todas as regiões do Brasil.
A Mediana do Brasil como um todo foi 16,7, a mesma obtida na região Norte.
A maior Mediana foi obtida na região Sul ( 23,3 ), e a menor, na região Sudeste ( 15,0 ).
O Gráfico 6.10 representa a distribuição das notas nas questões discursivas no Componente de Conhecimento Específico.
A moda dessa distribuição ocorre no intervalo [ 0 ; 10 ], com cerca de 35 \% do total de participantes.
A distribuição, \textbf{grosso} modo, é decrescente com a nota.
Observa -se que nesses intervalos estão computadas também as ocorrências de respostas em branco para as três questões discursivas de Conhecimentos Específicos.
A análise de cada uma destas questões será feita a seguir. 6.3.2.1 Análise da Questão Discursiva 3 do Componente de Conhecimento Específico Na questão 3, cujos resultados aferidos encontram -se descritos na Tabela 6.18, a Média dos estudantes de todo o Brasil foi 47,2, a questão de melhor desempenho dentre as três discursivas de Conhecimentos Específicos.
A menor Média nessa questão foi obtida pelos alunos da região Sudeste ( 43,2 ), enquanto a maior Média foi obtida na região Sul ( 55,9 ).
Quanto à variabilidade das notas, o Desvio padrão de todo o Brasil foi 34,8.
O maior Desvio padrão foi obtido na região Norte ( 35,8 ), enquanto o menor foi obtido na região Sul ( 33,6 ).
A nota Máxima, 100,0 pontos, foi alcançada por, pelo menos, um aluno de todas as regiões.
A Mediana do Brasil como um todo foi 50,0, a mesma para as regiões Norte, Nordeste e Centro-Oeste.
Nas regiões Sudeste e Sul, a Mediana foi, respectivamente, 40,0 e 65,0.
A nota Mínima ( 0,0 ) foi a mesma em todas as regiões do Brasil.
O Gráfico 6.11 mostra a distribuição das notas na questão discursiva 3, do Componente de Conhecimento Específico da área de Engenharia Civil.
A distribuição é multimodal, com moda principal no intervalo ( 90 ; 100 ], acompanhado de muito perto pelos intervalos ( 10 ; 20 ] e ( 60 ; 70 ], que são dois dos máximos locais.
Destaca -se, também, que cerca de 20 \% dos alunos ou deixaram a questão em branco ou obtiveram a nota zero. 6.3.2.2 Comentários sobre as respostas à Questão Discursiva 3 A Questão 3 abordou um tema de grande importância para o Engenheiro Civil, envolvendo um conteúdo básico, pertencente ao núcleo profissional de qualquer curso de Engenharia Civil, no país.
O enunciado da questão apresentou -se de forma clara, com informações suficientes para a formulação da resposta pelos estudantes.
Considera -se que esta foi uma questão de nível médio, pois solicitava duas informações para cada estado do concreto, fresco e endurecido, o nome e a descrição dos seus objetivos.
A questão se adequa às Diretrizes Curriculares do curso de Engenharia Civil.
A diversidade de respostas dos estudantes foi contemplada pelo padrão de respostas corretas possíveis.
Em sua maioria, os estudantes utilizaram uma linguagem adequada em relação aos termos específicos e técnicos e puderam expressar os conhecimentos adquiridos em relação ao tema abordado.
Muitos estudantes responderam à questão de forma correta e tiveram um bom desempenho.
Alguns cometeram alguns equívocos e, portanto, tiveram um desempenho Muitos estudantes responderam à questão de forma correta e tiveram um bom desempenho.
Alguns cometeram alguns equívocos e, portanto, tiveram um desempenho médio ou insuficiente.
A grande maioria apresentou uma descrição dos objetivos de forma clara e correta, no entanto, não citou os nomes dos respectivos ensaios.
Também foi identificado que muitos estudantes fizeram uma descrição do método dos ensaios, informação não solicitada na questão.
Houve quem citasse outros ensaios, que podem ser realizados no concreto, com finalidades diferentes da solicitada pela questão, que era o controle de qualidade.
Como se trata de uma questão discursiva, houve uma grande diversidade de respostas, em relação à forma de escrita, porém dentro dos aspectos exigidos pelo comando da questão e esperados pelo padrão de resposta, de forma parcial e/ou integral.
O maior percentual de respostas, superando 40 \%, foi de nível mediano, com notas maiores do que 30 e menores do que 70.
O percentual de respostas insuficientes e boas foi bem próximo, aproximadamente de 30 \% em cada grupo.
Um pouco mais do que 10 \% das respostas estavam completamente erradas e receberam nota zero.
Cabe ressaltar que nesta questão, também o índice de respostas em branco foi inferior a 9 \%, bem inferior ao das outras duas questões discursivas de Conhecimentos Específicos.
Apesar de a maioria das notas ter sido mediana, os estudantes demonstraram conhecimento sobre o assunto, mesmo que de forma parcial.
Alguns apresentaram uma descrição do método dos ensaios, informação não solicitada na questão, e outros fizeram uma descrição dos objetivos de forma clara e correta, apesar de, não citarem o nome dos respectivos ensaios. 6.3.2.3 Análise da Questão Discursiva 4 do Componente de Conhecimento Específico A Tabela 6.19 contém as informações relativas à questão 4 do conjunto de questões do Componente de Conhecimento Específico.
O desempenho dos estudantes nessa questão foi inferior ao da questão 3 e um pouco superior ao da questão 5.
A Média geral do Brasil foi 7,1, sendo a menor Média registrada na região Norte ( 3,9 ), e a maior, na região Nordeste ( 9,1 ).
A nota Máxima ( 100,0 ) foi atingida por, pelo menos, um concluinte de todas as regiões.
A Mediana em todo o Brasil foi zero, o mesmo valor foi obtido nas cinco Grandes Regiões, indicando que mais de 50 \% dos respondentes de todas as regiões tiraram nota zero.
Também foi zero a nota mínima de todas as regiões.
O Gráfico 6.12 representa a distribuição de notas da questão discursiva 4, do Componente de Conhecimento Específico.
Essa distribuição tem moda na classe de estudantes que, tendo respondido à questão, obtiveram nota zero.
Destaca -se o fato de que mais de 80 \% dos alunos ou deixaram a questão em branco ou obtiveram a nota zero. 6.3.2.4 Comentários sobre as respostas à Questão Discursiva 4 A questão abordava conceitos teóricos específicos, que são introduzidos no ciclo básico e ao longo de um curso de Engenharia Civil, pois tais conceitos são estudados em disciplinas de fenômenos de transporte ou correlatas, hidráulica e saneamento, sendo que, tais disciplinas pertencem ao Núcleo de Conteúdos Profissionalizantes, conforme RESOLUÇÃO CNE/CES 11, DE 11 DE MARÇO DE 2002, no item XXII - Hidráulica, Hidrologia Aplicada e Saneamento Básico.
Ao longo do curso, tais conceitos podem ser vistos de forma pontual em outras disciplinas, muitas vezes utilizando sua aplicabilidade em tabelas.
Muitos cursos de Engenharia no país são divididos em ênfases e, neste caso, o conteúdo avaliado pela Questão 4 pertence à ênfase de Recursos Hídricos.
Os conhecimentos avaliados pela questão fazem parte das Diretrizes Curriculares, porém era uma questão \textbf{numérica} e que dependia de um conhecimento muito específico na área de recursos hídricos, dentro do curso de Engenharia Civil, tornando a questão difícil.
O enunciado e a formulação da pergunta se apresentaram de forma clara e com uma linguagem adequada para o entendimento por parte do estudante.
Algumas informações extras ( relativas aos conceitos físicos das grandezas envolvidas ) poderiam ter sido indicadas no enunciado, o que possibilitaria que os estudantes tivessem maior êxito em suas respostas.
Com um padrão de respostas que contemplava as respostas possíveis para área molhada, perímetro molhado e raio hidráulico, esperava -se contemplar uma quantidade maior de respostas.
Entretanto, foi observado um elevado percentual de estudantes que deixou a questão em branco.
Muitos outros deixaram claro que não detinham o conhecimento das grandezas físicas e utilizaram fórmulas equivocadas, não chegando, portanto, à resposta correta, nem de forma parcial.
Os estudantes, em sua maioria, tiveram muita dificuldade em responder à questão, demonstrando pouco conhecimento em relação ao assunto abordado.
Muitos cometeram erros em relação às formulações das grandezas \textbf{numéricas}, seja para a área molhada, para o perímetro molhado ou para o raio hidráulico.
Merece destaque o elevado número de estudantes que sequer respondeu à questão ( superior a 35 \% ), o que talvez se explique por ser necessário o conhecimento das fórmulas para se obtenção da resposta correta.
Mais de 90 \% das respostas foram de nível insuficiente e, dentre essas, muitas ( mais do que 75 \% ) com nota zero.
O percentual de respostas medianas não superou 2 \%, e o percentual de respostas boas foi de aproximadamente 7 \%. 6.3.2.5 Análise da Questão Discursiva 5 do Componente de Conhecimento Específico A Tabela 6.20 contém as informações relativas à questão 5 do conjunto do Componente de Conhecimento Específico.
O desempenho dos estudantes de todo o Brasil nesta questão foi o mais baixo dentre as três questões discursivas desse componente.
A nota Média dos estudantes de todo o Brasil foi 2,2.
A maior Média foi registrada na região Nordeste ( 3,2 ), enquanto a menor Média foi registrada na região Sudeste ( 1,6 ).
Quanto à variabilidade das notas, o Desvio padrão dos alunos do Brasil, como um todo, foi 9,4.
Enquanto o maior desvio foi encontrado na região Nordeste ( 11,5 ), o menor foi encontrado na região Sudeste ( 7,9 ).
A Mediana para o Brasil e em todas as regiões foi zero, associada ao excessivo número de notas zero, como mostra o Gráfico 3.13.
Para o conjunto de alunos de Engenharia Civil do Brasil, a nota Máxima foi 100,0, e a nota Mínima ( 0,0 ).
Os mesmos limites foram encontrados em todas as regiões.
A Mediana também foi zero em todas as regiões, indicando que mais de metade dos alunos receberam essa nota, em cada região.
O Gráfico 6.13 representa a distribuição das notas da questão discursiva 5 do Componente de Conhecimento Específico.
Essa distribuição, como todas as relativas às questões discursivas dessa parte da prova, também tem moda principal nas notas zero.
Destaca -se que a frequência de notas zero somada a de estudantes que não responderam à questão chega a aproximadamente 90 \% do total de participantes. 6.3.2.6 Comentários sobre as respostas à Questão Discursiva 5 A Questão 5 abordava conceitos de ensaios de campo em solos.
A questão se adequa às Diretrizes Curriculares, cujos conceitos são estudados em disciplinas da área de Geotecnia, tais como Mecânica dos Solos, Obras de Terra e correlatas, as quais pertencem ao Núcleo de Conteúdos Profissionalizantes, conforme RESOLUÇÃO CNE/CES 11, DE 11 DE MARÇO DE 2002, no item XVII - Geotecnia, porém avalia conteúdo muito específico, dentro da área de transportes.
Tais ensaios costumam ser abordados em disciplinas do núcleo profissional básico, muitas vezes de forma superficial, em disciplinas de Mecânica dos solos.
Posteriormente, é comum ser retomado no ciclo profissional, na área de transportes.
Geralmente, esse conteúdo é abordado com mais profundidade no último ano do curso, em disciplinas de pavimentação ou correlatas.
No entanto, como se trata de assunto muito específico, muitos cursos no Brasil não contemplam essa abordagem de forma tão detalhada.
Além disso, pelo fato de muitos cursos serem divididos em ênfases e este conteúdo pertencer à área de transportes, há estudantes que não retomam, de forma mais aprofundada, conhecimentos específicos de ensaios de campo.
O enunciado da questão foi redigido em linguagem adequada para o entendimento do estudante.
Apesar disso, a maneira como a pergunta estava formulada parece ter confundido uma grande parte dos estudantes, uma vez que a questão solicitava que fosse explicado como deveria ser feito o controle em campo de três propriedades do solo.
Muitos estudantes responderam sobre o controle desses \textbf{parâmetros} em laboratório, outros descreveram os objetivos do controle dessas propriedades, e houve muitas respostas descrevendo o método de compactação dos solos.
A incompreensão do comando pode ter sido uma das razões do alto índice de respostas em branco e de respostas totalmente equivocadas.
Para os respondentes, esta foi uma questão de nível difícil, em comparação com ás demais questões de formação específica, em virtude das respostas identificadas.
Poucos estudantes utilizaram uma linguagem adequada em relação aos termos específicos e técnicos do problema proposto.
O padrão de respostas contemplou a diversidade de respostas dos estudantes, tornando possível diferenciar respostas insuficientes, medianas e boas, com pontuações adequadas para as diferentes respostas parcialmente corretas.
A maioria dos estudantes teve muita dificuldade em responder à questão, demonstrando pouco conhecimento em relação ao assunto abordado.
Observou -se, também, uma grande diversidade de respostas.
Como já mencionado, muitas respostas abordavam o controle dos \textbf{parâmetros} através de testes em laboratório, enquanto a questão solicitava respostas sobre os procedimentos de controle no campo.
Alguns descreveram, de forma equivocada, os objetivos do controle dessas propriedades para o controle tecnológico das obras de pavimentação.
Também houve muitas respostas descrevendo o método de compactação dos solos.
Um percentual pequeno de estudantes conseguiu responder à pergunta da questão de forma correta, parcial ou total.
Mais de 90 \% das respostas receberam notas baixas, sendo elevado o percentual de respostas com nota zero, acima de 55 \%, e também o de respostas em branco, superior a 35 \%.
Já o percentual de respostas medianas não superou 4 \%, e o de respostas boas foi menor do que 1 \%.
A Questão 5 foi a de desempenho mais baixo dentre as três questões discursivas de conhecimentos específicos do Enade 2017 de Engenharia Civil. 6.3.3 Considerações Finais As questões propostas para os formandos de Engenharia Civil no Enade 2017 mostraram quão importante é a clareza e a completude dos enunciados.
Muitos estudantes apresentaram respostas diversas do que se esperava nas três questões e, na Questão 4, talvez a apresentação de dados que normalmente seriam consultados por um profissional, pudesse ter levado a uma quantidade maior de respostas corretas.
Quanto à seleção dos conteúdos, foi observado que as questões abordaram temas bastante específicos que, apesar de fazerem parte da Formação Geral inicial, necessitariam de aprofundamento, ao longo da formação profissional, para que os estudantes se sentissem seguros para retomá -los em uma prova.
Também merece menção a dificuldade encontrada pelos estudantes na Questão 4, provavelmente, por ser uma questão \textbf{numérica}, que exigia a aplicação de fórmulas, mesmo que simples, envolvendo área e perímetro de círculos.
Além disso, conteúdos muito específicos, são problemáticos, tendo em vista a grande diversidade dos currículos dos cursos de Engenharia Civil no país.
Cabe destacar que alguns cursos têm foco em uma formação sólida, porém ampla, e envolvendo os principais conceitos ( cursos sem ênfase ), e outros, que são voltados para uma formação mais especializada ( cursos com ênfases ).
Essas considerações podem ajudar a explicar o melhor desempenho dos estudantes na Questão 3 e o baixíssimo desempenho revelado nas Questões 4 e 5 do presente exame.

% matched lemmas: grosso, numérico, parâmetro
\end{textsample}
